\begin{lemma}
Let $D>0$ be an integer, let $\mathcal H\in H(3,2,D)$ have finite vertex
and edge sets, and let $f$ be an edge. Suppose every vertex of $f$ has degree
$D$ and $f$ has exactly $3(D-1)$ neighboring edge indices. Put
$B=D+18D^{2/3}$ and $r=1+18D^{-1/3}$.
There are $n\ge3$ and a real $3$-by-$n$ matrix $A$ with the following
properties. There is an injective enumeration $v_1,v_2,v_3$ of $f$ and
slots $x_1,\ldots,x_n\in V\sqcup\{\bot\}$ such that no real vertex
occurs twice, and the real vertices occurring are exactly
\[
\{y\notin f:D^{2/3}\le\sum_{i=1}^3\deg_{\mathcal H}(v_i,y)\}.
\]
Here $\bot$ is an additional dummy symbol, not a vertex. The entries are
$A_{ij}=\deg_{\mathcal H}(v_i,x_j)/B$ for real slots and zero for dummy
slots. Each entry is in $[0,1]$, every column sum is at most one, the column
sums are nonincreasing, and the sum of all entries is at most three.
Writing
\[
Q=\sum_j\sum_{i<i'}A_{ij}A_{i'j},\quad
P=\operatorname{per}(A_{[3],[3]}),\quad
L=\sum_{i=1}^3\sum_{j=1}^3 A_{ij},
\]
one has
\[
b_{L(\mathcal H),3D}(f)\le
1-\frac29+\frac29D^{-1/3}+\frac{8}{9D}
 +\frac{r^2Q}{9}+\frac{r^3P}{27}-\frac{rL}{27}.
\]
\end{lemma}
\begin{proof}
Enumerate the three distinct vertices of $f$, using its uniformity from
\noderef{HypergraphClass}. Let $a_i(y)=\deg_{\mathcal H}(v_i,y)$ and
let $X$ consist of vertices outside $f$ lying in a neighboring edge.
By \noderef{SaturatedEdgeStructure}, each neighboring edge meets $f$ once,
each column sum on $X$ is at most $D$, and
$\sum_{y\in X}a_i(y)=2(D-1)$. Codegrees are nonnegative counts.
Put $X_s=\{y\in X:\sum_i a_i(y)<D^{2/3}\}$ and $X_b=X\setminus X_s$.
Since $D^{2/3}>0$, a vertex outside $f$ satisfying the reverse threshold
has a positive codegree with some $v_i$. An edge counted by that codegree
is different from $f$ and meets $f$, so this vertex is in $X$.
Consequently $X_b$ is exactly the set displayed in the statement.

Let $m=|X_b|$. Enumerate this finite set by repeatedly choosing a vertex of
largest remaining column sum. At each of the $m$ steps the remaining
nonempty finite set has a maximum; deleting the chosen vertex gives a
list without repetitions and the list exhausts $X_b$. Its column sums
are nonincreasing. Set $n=\max(3,m)$, use these $m$ vertices as the first
slots, and put the dummy symbol in all other slots. Thus the slots are
literally Option-valued, and a dummy slot contributes zero rather than
the codegree of an arbitrarily chosen vertex. Nonnegative real column sums
remain nonincreasing after the zeros are appended.

Define $d_{ij}=a_i(x_j)$ on real slots and zero on dummy slots, and set
$A_{ij}=d_{ij}/B$. Since $B\ge D>0$, the column bound gives
$0\le A_{ij}\le1$ and column sums at most one. Sorting is preserved by
division by $B$. We verify the big-column row estimate directly for all
positive $D$. Summing the threshold over $X_b$ and then using the three
row identities gives $|X_b|D^{2/3}\le6(D-1)\le6D$, so
$|X_b|\le6D^{1/3}$. Fix $v_i$. Among its $D-1$ neighboring edges, an
edge with two vertices of $X_b$ is specified by an unordered pair from
$X_b$. There is at most one edge index per pair: two would share the
three vertices consisting of that pair and $v_i$, violating the
$2$-simplicity in \noderef{HypergraphClass}. The number $q_i$ of such
edges is therefore at most $\binom{|X_b|}{2}\le |X_b|^2/2\le18D^{2/3}$.
Every other edge through $v_i$ contains at most one vertex of $X_b$.
Each codegree incidence in $X_b$ comes from one of these $D-1$ edges,
because a member of $X_b$ is outside $f$. Thus
$\sum_{y\in X_b}a_i(y)\le(D-1)+q_i\le D+18D^{2/3}=B$.
Each real vertex occurs once and dummy slots contribute zero, so every
row of $A$ sums to at most one. Summing the three rows bounds its total
sum by three.

Write $P_f,T_f$ for the independent pair and triple counts in the
line-graph neighborhood. For a subset $Y$ of $X$, put
$w(Y)=\sum_{y\in Y}\sum_{i<i'}a_i(y)a_{i'}(y)$.
The unordered vertex pairs in the statement of
\noderef{SaturatedLineGraphPairBound} are exactly the three sets
$\{v_i,v_{i'}\}$ with $i<i'$, each occurring once, since the enumeration
of $f$ is injective and surjective. Its outside-vertex set is $X$: an
edge containing a vertex outside $f$ cannot itself be $f$.
Its integer inequality therefore reads
$P_f\ge3D^2-6D+3-w(X)$, also as a real inequality.
For completeness the small-column estimate also holds for every positive
$D$: for three real numbers $a,b,c$, the sum of the three squares
$(a-b)^2+(a-c)^2+(b-c)^2\ge0$ gives
$ab+ac+bc\le(a+b+c)^2/3$. A small column has nonnegative sum
$s<D^{2/3}$, so its pair weight is at most $D^{2/3}s/3$.
Summing over $X_s$, the total of these column sums is at most
$6(D-1)\le6D$. Hence $w(X_s)\le2D^{5/3}$.
Since $X_s$ and $X_b$ partition $X$, $w(X)=w(X_s)+w(X_b)$, and the
pair inequality above yields
\[
P_f\ge3D^2-6D+3-w(X_b)-2D^{5/3}.
\]

Select the first three slots. Each real selected slot is in $X_b$ and is
unique. If a selected slot is dummy then $m<3$, so all real vertices have
already been selected. The first three column sums decrease, and any
unselected column is at most the third by the sorting construction.
All hypotheses of \noderef{ThreeUniformTwoSimpleTripleBound} now hold,
including the unique root intersection, column bounds and row identities
supplied above. Its conclusion is
\[
T_f\le\operatorname{per}(d_{[3],[3]})+3D^3+6D^2
       -D^2\sum_{i,j\le3}d_{ij}.
\]
The permutation sum there is the permanent: a product over the three
row indices is precisely its three displayed factors. This also holds
with dummy slots, whose entries are zero.

The formula in \noderef{LocalBParameter} and the degree $3(D-1)$ give
$b=3(D-1)/(3D)-P_f/(3D)^2+T_f/(3D)^3$.
Use the pair lower bound and triple upper bound, then replace the degree
term by its upper bound $1$ and drop the negative term $-3/(9D^2)$.
This yields the consistent looser bound with constant part
\[
1-\frac{3D^2}{9D^2}
 +\frac{2D^{5/3}+6D}{9D^2}
 +\frac{3D^3+6D^2}{27D^3}
=1-\frac29+\frac29D^{-1/3}+\frac8{9D}.
\]
The remaining terms are $w(X_b)/(9D^2)$,
$\operatorname{per}(d_{[3],[3]})/(27D^3)$, and
$-\sum_{i,j\le3}d_{ij}/(27D)$. Each pair product has two factors and
each permanent product has three, so $d=BA$ gives respectively
$B^2Q$, $B^3P$, and $BL$ in their numerators. Finally
$B/D=1+18D^{-1/3}=r$, by the exponent laws for positive $D$.
These substitutions give exactly the asserted inequality.
\end{proof}
